En un experiment d'efecte fotoelèctric, fem incidir una llum monocromàtica sobre el càtode d'una cèl·lula fotoelèctrica. Quan s'il·lumina el càtode amb llum blava de longitud d'ona 460~nm, s'observa un corrent fotoelèctric. Aplicant una diferència de potencial de com a mínim 0,50~V entre càtode i ànode, el corrent desapareix completament.

\begin{apartat}{1,25}
Calculeu l'energia mínima necessària per a arrencar un electró del càtode (funció de treball) i la velocitat màxima amb què surten els electrons emesos quan no s'aplica cap diferència de potencial.
\end{apartat}

\begin{apartat}{1,25}
Determineu quina diferència de potencial caldria aplicar per a anul·lar el corrent fotoelèctric si primer el càtode s'il·lumina amb llum verda de longitud d'ona 520~nm, i després amb llum vermella de longitud d'ona 650~nm. Justifiqueu els resultats obtinguts.
\end{apartat}

\dades{%
  $|e| = 1,602\times 10^{-19}\ \mathrm{C}$.\\
  $m_\mathrm{e} = 9,11\times 10^{-31}\ \mathrm{kg}$.\\
  $c = 3,00\times 10^{8}\ \mathrm{m\,s^{-1}}$.\\
  $h = 6,63\times 10^{-34}\ \mathrm{J\,s}$.}
